\section{Exercise 2}

Mảng khai báo bằng \texttt{.word 0:5} gồm 5 phần tử, mỗi phần tử 4 byte. Với chỉ số $i$ lưu trong \texttt{\$t0}, địa chỉ phần tử thứ $i$ là $\text{base} + 4i$, với phép nhân 4 được thay bằng dịch trái 2 bit:

\begin{lstlisting}[language=MIPS,numbers=none]
  sll $t2, $t0, 2   # $t2 = 4i
  add $t2, $s0, $t2 # $t2 = base + 4i
  sw $t1, 0($t2)    # cất $t1 vào phần tử thứ i
\end{lstlisting}

Mỗi lời nhắc gồm 3 syscall: syscall 4 in \texttt{"Please input element "}, syscall 1 in chỉ số, syscall 11 in dấu \texttt{":"}:

\begin{lstlisting}[language=MIPS,numbers=none]
input_element_loop:
  li $v0, 4            # syscall 4: in chuỗi
  la $a0, input_prompt # "Please input element "
  syscall

  li $v0, 1            # syscall 1: in số nguyên
  move $a0, $t0        # in chỉ số đang duyệt
  syscall

  li $v0, 11           # syscall 11: in ký tự
  li $a0, ':'          # in dấu hai chấm
  syscall
\end{lstlisting}

Chỉ số nhập vào được dùng trực tiếp để tính offset rồi nạp phần tử:

\begin{lstlisting}[language=MIPS,numbers=none]
  li $v0, 5
  syscall

  move $t0, $v0     # đưa chỉ số vào $t0
  sll $t2, $t0, 2   # dịch trái 2 bit để nhân 4
  add $t3, $s0, $t2 # tính địa chỉ phần tử

  li $v0, 1
  lw $a0, 0($t3)
  syscall
\end{lstlisting}

Với mảng $5, 10, 15, 20, 25$, nhập chỉ số 1 rồi nhập chỉ số 4:

\begin{lstlisting}[language={},numbers=none]
Please input element 0:5
Please input element 1:10
Please input element 2:15
Please input element 3:20
Please input element 4:25
Please enter index:1
10
\end{lstlisting}

\begin{lstlisting}[language={},numbers=none]
Please input element 0:5
Please input element 1:10
Please input element 2:15
Please input element 3:20
Please input element 4:25
Please enter index:4
25
\end{lstlisting}
